\section{Loop 3: Session Reflection}

\label{sec:loop3}

\subsection{Motivation and Design}

Session reflection occupies the final 2 minutes of each session. Its purpose
is to synthesize the session's telemetry and context into a compact structured
record that Loop~4 can aggregate across sessions. The key design constraint is
that reflection must use \emph{constrained templates} rather than open-ended
prompts.

Open-ended reflection prompts (``What did you learn this session?'') produce
outputs that are useful as prose but difficult to aggregate. Two sessions with
similar issues will produce different prose descriptions of those issues,
making cross-session pattern detection unreliable. Constrained templates with
fixed output schemas solve this.

\subsection{Reflection Template}

The session reflection uses the following three-question template, which
the LLM answers with structured JSON output:

\vspace{0.5em}
\noindent\textbf{Question 1: What worked?}
Enumerate tool invocations or strategies with outcome $=$ \textsc{success}
that were reused $\geq 2$ times in the session. Report as:
$\{$\texttt{tool}, \texttt{strategy}, \texttt{success\_count}, \texttt{why\_effective}$\}$.

\noindent\textbf{Question 2: What failed repeatedly?}
Enumerate $(t, f)$ pairs with count $\geq 2$ that did not trigger the BIN
loop (i.e., threshold not yet reached). Report as:
$\{$\texttt{tool}, \texttt{failure\_mode}, \texttt{count}, \texttt{hypothesis}$\}$.

\noindent\textbf{Question 3: What should change?}
Based on the above, enumerate proposed changes to the tool library, system
prompt, or behavior patterns. Each proposal must reference specific evidence
from Questions 1 and 2. No proposal without evidence is recorded.

\vspace{0.5em}

\begin{definition}[SessionReflection Record]
\label{def:session_reflection}
A \textsc{SessionReflection} record $\rho$ is a tuple:
\begin{equation}
\rho = \bigl( \text{session\_id},\; \text{successes},\; \text{failures},\;
              \text{proposals},\; \text{telemetry\_digest} \bigr),
\end{equation}
where $\text{telemetry\_digest}$ is a compact statistical summary of
$\mathcal{T}^{(s)}_{\text{cur}}$ (tool usage histogram, failure rate per tool,
session duration, task count) that is appended to the reflection record for
Loop~4 consumption without requiring Loop~4 to replay the full log.
\end{definition}

\subsection{Constrained Templates vs.\ Open-Ended Reflection}

\begin{theorem}[Template Aggregation Advantage]
\label{thm:template_aggregation}
Let $R_T$ be the set of reflections produced by the constrained template, and
$R_O$ be the set of reflections produced by an open-ended prompt, over $N$
sessions with the same underlying pattern distribution. Then the expected
precision of cross-session pattern detection is:
\begin{equation}
P(\text{detect} \mid R_T) \geq P(\text{detect} \mid R_O) + \Delta,
\end{equation}
where $\Delta > 0$ depends on the entropy of the natural language description
of patterns and the schema alignment between sessions.
\end{theorem}

\begin{proof}[Proof sketch]
Cross-session pattern detection requires matching descriptions of similar events
across sessions. For templated reflections, matching reduces to structured key
equality (same tool name, same failure mode code). For open-ended reflections,
matching requires semantic similarity, which has non-trivial false-negative
rate $\varepsilon > 0$ (e.g., ``permission denied on file writes'' and
``cannot write to read-only path'' describe the same pattern but are
lexically dissimilar). Therefore $P(\text{detect} \mid R_T) = 1 - O(|\mathcal{V}_T|^{-1})$
while $P(\text{detect} \mid R_O) \leq 1 - \varepsilon$ for some $\varepsilon > 0$
that depends on the natural language variability of failure descriptions.
\end{proof}

% ============================================================================
% 8. LOOP 4: MAINTENANCE PASS
% ============================================================================
